% =============================================================================
\section{Introduction}
\label{sec:intro}
% =============================================================================
\cgal{} (Computational Geometry Algorithm Library) is an open source software project that provides access to efficient and reliable implementations of geometric algorithms in the form of a C++ library~\cite{cgal:eb-21}. \cgal{} has evolved through the years and now represents the state of the art in  computational geometry software. The software modules of \cgal{} rigorously adhere to the generic programming paradigm---a discipline that consists of the gradual lifting of concrete algorithms abstracting over details, while retaining the algorithm semantics and efficiency. The software modules of \cgal{} also follow the exact geometric-computation paradigm, which simply amounts to ensuring that errors in predicate evaluations do not occur; it guarantees robustness of the applied algorithms.  As a result, the software of \cgal{} is stable despite the use of (inexact) floating point arithmetic, found in standard hardware, it is complete as it handles all degenerate cases, which typically are in abundance, and it is efficient---all at the same time.

The \cgal{} \emph{arrangements} packages constitute a large component of the \cgal{} library---a component that is particularly intricate, partly due to the interplay between combinatorial algorithms and algebra~\cite{fhw-caass-12}. Arrangements are space subdivisions induced by curves and surfaces, which have been intensively studied in discrete and computational geometry~\cite{hs-a-18}, and have applications in various domains, from robotics~\cite{my-hs-atarr-95} and assembly planning~\cite{hlw-gfapm-00} through Geographic Information Systems (GIS)~\cite{DBLP:conf/gis/KreveldSW04} to protein structure determination~\cite{DBLP:journals/jcb/MartinYBZD11}, to mention just a few uses. The arrangements packages of \cgal{} have been developed since the early days of \cgal, first for planar arrangements and maps~\cite{DBLP:journals/jea/FlatoHHNE00, DBLP:journals/comgeo/WeinFZH07}, Boolean operations, and Minkowski sums~\cite{bfhhm-epsph-15}. Then, envelopes of surfaces in three-dimensions have been added~\cite{m-rgece-06}. Finally, a major effort has been undertaken to supports two-dimensional arrangements on (not necessarily planar) surfaces~\cite{bfhks-apsca-10,bfhmw-apsgf-10}.

% There is a large community of users of the \cgal{} arrangements packages in academia and industry, and these packages seem to be helpful in many diverse projects. 
There is a large community of users of the \cgal{} arrangements packages in academia and industry engaged in many diverse projects. 
In all the projects that we are aware of, the packages have been used by C++ experts or by users who have been strongly supported by such experts from the \cgal{} developers community, including GeometryFactory.\footnote{GeometryFactory is the company established to market and support the usage of \cgal{} packages.} 
% In our experience, incorporating \cgal{} in education or in academic projects lacking C++ experts has been rugged. We attribute the difficulties that arise in such settings to the highly sophisticated usage of C++ in \cgal.
In our experience, incorporating \cgal{} in education or in academic projects lacking C++ experts has been rugged due to the highly sophisticated usage of C++ in \cgal.
The binding project that we describe in this paper is meant to resolve these difficulties and make the arrangements packages more accessible to a wider audience of programmers. We already see, and report below on, first signs that the project is on a good trajectory toward achieving this goal.

% -----------------------------------------------------------------------------
\subsection{C++ Python Bindings}
\label{ssec:intro:bind}
% -----------------------------------------------------------------------------
Python, like many other languages, was originally defined by a reference implementation. The original reference implementation of the Python programming language is CPython. To date, CPython is still the default and commonly used reference implementation of the language, although new rivals are emerging, see, e.g., PyPy.\footnote{PyPy is a replacement for CPython; see \url{https://www.pypy.org}.} As the Python language evolved, its definition has tightened and a well defined specification has materialized.

Bindings are essentially wrapper libraries that bridge two programming languages, so that a software module developed in one language can be used in code written in another language, exploiting (i)~the strengths of both languages, and (ii) the availability of modules developed in the former language. C++ Python bindings enable (i) the invocation of C++ functions, and (ii)~the access to C++ variables, from Python code taking advantage of Python's quick development cycles and C++'s high performance.

Python is sufficiently efficient for many tasks; nevertheless low-level code written in Python tends to be too slow, largely due to Python's extremely dynamic nature. In particular, low-level computational loops are simply infeasible. In addition, the magnitude of existing and well-tested code in Fortran, C, or C++, is enormous. The Python/C API, the module that enables the usage of external modules developed in Fortran, C, or C++, from Python code, or embedding Python code in code written in other languages, is very rich.  Indeed, the Python/C API module of CPython exposes a vast amount of details of CPython internals. It enables the efficient exploitation of existing code in Fortran, C, or C++, 
% as well as the replacement of critical sections written in Python to be written in these languages when speed is essential. 
as well as the replacement of critical sections written in Python when speed is essential.
Wrapping existing code has traditionally been the domain of Python experts due to the steep learning curve, which characterizes its Python/C API. Although using such wrappers is possible without ever knowing their internals, providing such wrappers creates a sharp line between developers using Python and developers using C or C++ with the Python/C API. Several tools that facilitate the creation of such wrappers have been introduced; The three most propitious tools are listed In Section~\ref{sec:comp}; additional tools are listed in Section~\ref{sec:tools}.

When running Python code that uses bindings for some C++ modules, one or more libraries that provide the bindings must be accessible.  A software module in C++ that adheres to the generic programming paradigm consists of function and class templates; these templates are instantiated at compile time of the binding libraries. In other words, the types of the C++ objects that are bound, that is, instances (instantiated types) of C++ function and class templates, must be known when the bindings are generated.

% -----------------------------------------------------------------------------
\subsection{Bindings for \cgal}
\label{ssec:intro:cgal}
% -----------------------------------------------------------------------------
Instantiated types in \cgal{} are characterized by long instantiation chains of C++ class templates. An instantiated type is a template that has one or more template parameters and every parameter is substituted by another type that is typically an instance of another template. While the number of models of most concepts is small, the number of potential types of objects that could be bound, and thus must be supported by the bindings, is enormous. Offering bindings for all these types in advance is practically impossible. Moreover, in some cases models need to be extended with types provided by the user. Naturally, bindings for user-extended models must be generated dynamically.

\cgal{} uses \cmake~\cite{mh-mc-08} for the build process.  Since version~5.0, \cgal{} is header-only by default, which means that an application that depends on \cgal{} only depends on (i) the source code of \cgal{} that resides in header file,
% (the name of which typically end with the suffix \myLstinline{.hpp}),
and (ii) native configuration files generated by \cmake{} (in particular, the execution of \myLstinline{cmake}). In other words, there is no need to compile any library before an application that depends on \cgal{} and written in C++ is built.\footnote{Some dependencies of \cgal{} might need to be installed in advance.} On the other hand, when running an application written in Python that depends on \cgal{} the libraries that contain the bindings must be accessible. Typically, those bindings are generated ahead of time with little knowledge about the application itself. This would impose a severe limitation in cases like ours, where the number of objects to be bound is large. In cases where the type of the objects to be bound are closely tied with the application itself, generating bindings ahead of time is, naturally, impossible. Our system enables easy and convenient (re)generation of bindings, thus alleviating the burden caused by these problems.

% =============================================================================
\subsection{Contribution}
\label{ssec:intro:contribution}
% =============================================================================
We introduce C++ Python bindings that enable the convenient, efficient, and reliable use of software modules from the \cgal{} \iiDArrangementsPackage{}, \iiDRegularizedBooleanSetOperationsPackage{}, and \iiDMinkowskiSumsPackage{} packages (and additional dependant packages) in Python code. There are different tools that facilitate the creation of C++ Python bindings. We present a short study that compares three main tools, which leads to the method of choice.
%
The binding themselves are implemented in C++ and their implementation exploits advanced features in C++, which we explicate. This results with elegant and compact code that is easy to maintain and extend with additional bindings.
%
We present a system that dynamically generates bindings of desired objects according to user prescriptions, which enables the convenient use of any subset of bound object types concurrently.  With our system it is possible to generate a single library that contains bindings for instances of different \cgal{} templates,
% , e.g., an instance of the 2D arrangement class template and an instance of the 2D polygon-set class template, or several libraries when one is insufficient. If several instances of the same class template must be bound, several corresponding libraries must be generated, one for each instance.
or several libraries (which can be used concurrently), such that distinct libraries contain different instances of the same template.

% -----------------------------------------------------------------------------
\subsection{Outline}
\label{ssec:intro:outline}
% -----------------------------------------------------------------------------
The rest of this paper is organized as follows. A study that compares three tools for implementing C++ Python bindings for \cgal{} is presented in Section~\ref{sec:comp}. A sample of highlighted modern C++ techniques and idioms used in the binding implementation of several geometric procedures employed by the arrangements packages are described in Section~\ref{sec:code}. Applications of the new tools, which are geared toward making the CGAL code more accessible are presented in
Section~\ref{sec:app}. General instructions for dynamically generating bindings (see Section~\ref{sec:gen}), the description of the bindings generated by our system for various modules (see Section~\ref{sec:modules}), and the tight coupling between concepts and binding generation (see Section~\ref{sec:concepts}) are deferred to the appendix due to lack of space even though the content of these sections is essential for those who wish to study the entire work.

\begin{comment}
% -----------------------------------------------------------------------------
\subsection{Conventions}
\label{ssec:intro:conventions}
% -----------------------------------------------------------------------------
The paper is packed with code samples written either in the \cmake{}, Python, or C++ languages. We use several conventions to improve readability.  Code excerpt written in C++ is always numbered, while short code excerpts written in Python is sometimes preceded with \pyLstinline{>>>}. The C++ identifier \cppLstinline{pb} is an alias for the namespace \cppLstinline{boost::python}.
\end{comment}